\uuid{qL7Q}
\exo7id{2734}
\titre{exo7 2734}
\auteur{tumpach}
\organisation{exo7}
\datecreate{2009-10-25}
\isIndication{false}
\isCorrection{false}
\chapitre{Déterminant, système linéaire}
\sousChapitre{Système linéaire, rang}
\cpge{mpsi-mp2i.calcul-algebrique-trigonometrie.b, pcsi.calcul-algebrique-trigonometrie.b, ptsi.calcul-algebrique-trigonometrie.b, mpsi-mp2i.calcul-matriciel-systemes.c, pcsi.calcul-matriciel-systemes-lineaires.c, ptsi.calcul-matriciel-systemes-lineaires.c, mpsi-mp2i.matrices-applications-lineaires.c, pcsi.matrices-applications-lineaires.d, ptsi.matrices-applications-lineaires.d, bcpst1.systemes-lineaires.a, bcpst1.systemes-lineaires.c}
\module{Algèbre}
\niveau{L2}
\difficulte{}

\contenu{
\question{
R\'esoudre le syst\`eme~:
$$
\left(\begin{array}{ccccc}
1 & 0 & 2 & 1 & 1\\2 & 1 & 3 & -1 & 2\\ -2 & -1 & 1 & -3& 2\\3 & 2 & 0 & 1&-1
\end{array}
\right)
\left(\begin{array}{c}x \\ y \\ z \\ t\\u\end{array}\right) = \left(\begin{array}{c} 2 \\ 0 \\ 1 \\ 1 \end{array}\right).
$$
}
}
